%%%%%%%%%%%%%%%%%
% This is an sample CV template created using altacv.cls
% (v1.3, 10 May 2020) written by LianTze Lim (liantze@gmail.com). Now compiles with pdfLaTeX, XeLaTeX and LuaLaTeX.
% This fork/modified version has been made by Nicolás Omar González Passerino (nicolas.passerino@gmail.com, 15 Oct 2020)
%
%% It may be distributed and/or modified under the
%% conditions of the LaTeX Project Public License, either version 1.3
%% of this license or (at your option) any later version.
%% The latest version of this license is in
%%    http://www.latex-project.org/lppl.txt
%% and version 1.3 or later is part of all distributions of LaTeX
%% version 2003/12/01 or later.
%%%%%%%%%%%%%%%%

%% If you need to pass whatever options to xcolor
\PassOptionsToPackage{dvipsnames}{xcolor}

%% If you are using \orcid or academicons
%% icons, make sure you have the academicons
%% option here, and compile with XeLaTeX
%% or LuaLaTeX.
% \documentclass[10pt,a4paper,academicons]{altacv}

%% Use the "normalphoto" option if you want a normal photo instead of cropped to a circle
% \documentclass[10pt,a4paper,normalphoto]{altacv}

\documentclass[10pt,a4paper,ragged2e,withhyper]{altacv}

%% AltaCV uses the fontawesome5 and academicons fonts
%% and packages.
%% See http://texdoc.net/pkg/fontawesome5 and http://texdoc.net/pkg/academicons for full list of symbols. You MUST compile with XeLaTeX or LuaLaTeX if you want to use academicons.

% Change the page layout if you need to
\geometry{left=1.2cm,right=1.2cm,top=1cm,bottom=1cm,columnsep=0.75cm}

% The paracol package lets you typeset columns of text in parallel
\usepackage{paracol}
\usepackage{ragged2e}
\usepackage{hyperref}
\usepackage{tabto}

% Change the font if you want to, depending on whether
% you're using pdflatex or xelatex/lualatex
\ifxetexorluatex
  % If using xelatex or lualatex:
  \setmainfont{Roboto Slab}
  \setsansfont{Lato}
  \renewcommand{\familydefault}{\sfdefault}
\else
  % If using pdflatex:
  \usepackage[rm]{roboto}
  \usepackage[defaultsans]{lato}
  % \usepackage{sourcesanspro}
  \renewcommand{\familydefault}{\sfdefault}
\fi

% ----- LIGHT MODE -----
\definecolor{SlateGrey}{HTML}{2E2E2E}
\definecolor{LightGrey}{HTML}{666666}
\definecolor{PrimaryColor}{HTML}{001F5A}
\definecolor{SecondaryColor}{HTML}{0039AC}
\definecolor{ThirdColor}{HTML}{F3890B}
\definecolor{BackgroundColor}{HTML}{E2E2E2}
\colorlet{name}{PrimaryColor}
\colorlet{tagline}{PrimaryColor}
\colorlet{heading}{PrimaryColor}
\colorlet{headingrule}{ThirdColor}
\colorlet{subheading}{SecondaryColor}
\colorlet{accent}{SecondaryColor}
\colorlet{emphasis}{SlateGrey}
\colorlet{body}{LightGrey}
\pagecolor{BackgroundColor}   
% ----- DARK MODE -----
%\definecolor{BackgroundColor}{HTML}{242424}
%\definecolor{SlateGrey}{HTML}{6F6F6F}
%\definecolor{LightGrey}{HTML}{ABABAB}
%\definecolor{PrimaryColor}{HTML}{3F7FFF}
%\colorlet{name}{PrimaryColor}
%\colorlet{tagline}{PrimaryColor}
%\colorlet{heading}{PrimaryColor}
%\colorlet{headingrule}{PrimaryColor}
%\colorlet{subheading}{PrimaryColor}
%\colorlet{accent}{PrimaryColor}
%\colorlet{emphasis}{LightGrey}
%\colorlet{body}{LightGrey}
%\pagecolor{BackgroundColor}

% Change some fonts, if necessary
\renewcommand{\namefont}{\Huge\rmfamily\bfseries}
\renewcommand{\personalinfofont}{\small\bfseries}
\renewcommand{\cvsectionfont}{\large\rmfamily\bfseries}
\renewcommand{\cvsubsectionfont}{\large\bfseries}

% Change the bullets for itemize and rating marker
% for \cvskill if you want to
\renewcommand{\itemmarker}{{\small\textbullet}}
\renewcommand{\ratingmarker}{\faCircle}

%% sample.bib contains your publications
%% \addbibresource{sample.bib}

\begin{document}
    \name{Zhang Jiannan}
    \tagline{PhD Candidate}
    %% You can add multiple photos on the left or right
    \photoL{4cm}{JZhang}
    
    \personalinfo{
        \email{jiannan.zhang@tum.de}\smallskip
        \phone{+49-176-630-81506}
        \location{Munich, Germany}\\
        \linkedin{jiannan-zhang}
        %\github{johnDoe}
        %\dev{johnDoe}
        %\homepage{nicolasomar.me}
        %\medium{nicolasomar}
        %% You MUST add the academicons option to \documentclass, then compile with LuaLaTeX or XeLaTeX, if you want to use \orcid or other academicons commands.
        % \orcid{0000-0000-0000-0000}
        %% You can add your own arbtrary detail with
        %% \printinfo{symbol}{detail}[optional hyperlink prefix]
        % \printinfo{\faPaw}{Hey ho!}[https://example.com/]
        %% Or you can declare your own field with
        %% \NewInfoFiled{fieldname}{symbol}[optional hyperlink prefix] and use it:
        % \NewInfoField{gitlab}{\faGitlab}[https://gitlab.com/]
        % \gitlab{your_id}
    }
    
    \makecvheader
    %% Depending on your tastes, you may want to make fonts of itemize environments slightly smaller
    % \AtBeginEnvironment{itemize}{\small}
    
    %% Set the left/right column width ratio to 6:4.
    \columnratio{0.25}

    % Start a 2-column paracol. Both the left and right columns will automatically
    % break across pages if things get too long.
    \begin{paracol}{2}
        % ----- STRENGTHS -----
        \cvsection{\large Strengths}
            \cvtag{Modelling and Simulation}
            \cvtag{Flight Dynamics and Control}
            \cvtag{Control Theory}\\
            \cvtag{\small Control Allocation}\\
            \cvtag{MBSE}
            % \cvtag{Safety Critical Development}\\
            % \cvtag{Object Oriented Programming}
            % \cvtag{\small Nonlinear System Trim Linearization}\\
            
            % \medskip
            
            \cvtag{Flight Testing and Operation}
            \cvtag{UAV Piloting}
            
        % ----- STRENGTHS -----
        
        % ----- LEARNING -----
        \cvsection{\large Software Skills}
            \cvtag{MATLAB/Simulink}\\
            \cvtag{\normalsize Simulink Model Advisor}
            \cvtag{\normalsize Control Design Toolbox}
            \cvtag{\normalsize Math Symbolic Toolbox}
            \cvtag{Git}
            \cvtag{GitLab}
            \cvtag{Polarion}
            \cvtag{C/C++}
            \cvtag{Eclipse}
            \cvtag{CMake}
            \cvtag{Visual Studio Code}
            \cvtag{Python}
            \cvtag{Solidworks}
        % ----- LEARNING -----
        
        % ----- LANGUAGES -----
        \cvsection{\large Languages}
            \cvlang{English}{\small Professional Working}\\
            % \divider
            \medskip
            
            \cvlang{Chinese}{\small Native}\\
            \medskip
            % \divider
            
            
            \cvlang{German}{\small Goethe-Zertif. B1}
            \smallskip
            %% Yeah I didn't spend too much time making all the
            %% spacing consistent... sorry. Use \smallskip, \medskip,
            %% \bigskip, \vpsace etc to make ajustments.
            
        % ----- LANGUAGES -----
            
        % ----- REFERENCES -----
        % \cvsection{References}
        %     \cvref{Ref 1}{ref-1}
        %     \divider
            
        %     \cvref{Ref 2}{ref-2}
        %     \divider
            
        %     \cvref{Ref 3}{ref-3}
        %     \smallskip
        % ----- REFERENCES -----
        
        % ----- MOST PROUD -----
        % \cvsection{Most Proud of}
        
        % \cvachievement{\faTrophy}{Fantastic Achievement}{and some details about it}\\
        % \divider
        % \cvachievement{\faHeartbeat}{Another achievement}{more details about it of course}\\
        % \divider
        % \cvachievement{\faHeartbeat}{Another achievement}{more details about it of course}
        % ----- MOST PROUD -----
        
        % \cvsection{A Day of My Life}
        
        % Adapted from @Jake's answer from http://tex.stackexchange.com/a/82729/226
        % \wheelchart{outer radius}{inner radius}{
        % comma-separated list of value/text width/color/detail}
        % \wheelchart{1.5cm}{0.5cm}{%
        %   6/8em/accent!30/{Sleep,\\beautiful sleep},
        %   3/8em/accent!40/Hopeful novelist by night,
        %   8/8em/accent!60/Daytime job,
        %   2/10em/accent/Sports and relaxation,
        %   5/6em/accent!20/Spending time with family
        % }
        
        % use ONLY \newpage if you want to force a page break for
        % ONLY the current column
        \newpage
        
        %% Switch to the right column. This will now automatically move to the second
        %% page if the content is too long.
        \switchcolumn
        
        % ----- Objective -----
        \cvsection{Objective}
            \begin{quote} \justifying %\normalsize
                % Lilium:
                % With years of experience in flight control system development in both industry and academy, I am motivated to look for challenging opportunities where both my mindset for innovation and sound skills for safety critical system development can contribute, to exploit the promising future of urban air mobility.
                % APSIDE:
                With years of experience in the field of flight dynamics modelling, control and simulation in both industry and academy, I am motivated to look for challenging opportunities where both my mindset for innovation and sound skills for safety critical system development can contribute, to exploit the promising future of aerospace and air mobility industry.
            \end{quote}
        % ----- Objective -----
        
        % ----- EXPERIENCE -----
        \cvsection{Professional Experience}
            \cvevent{Research Associate }{| Institute of Flight System Dynamics}{06.2016 -- Now}{Munich, Germany}
            \begin{itemize} \justifying
                \item Participate in flight control system development for UAV prototypes.
                \item Lead of control allocation development for both general and safety critical systems.
                \item Organize and participate flight tests for small UAVs and work as test pilots.
                \item Support on building system modelling and  analysis tool chain.
            \end{itemize}
            \bigskip
            
            \cvevent{Intern Software Development}{| Phoenix-Wings GmbH}{06.2020 -- Now}{Ismaning, Germany}
            \begin{itemize}
                \item Focus on building a linearization tool in C++ to linearize an eVTOL dynamic model.
                \item Support in testing the software tool chains developed in the company including model-in-the-loop and software-in-the-loop simulation and ground control system.
                \item Attend flight tests and support in-flight ground control software testing.
                \item Support in preparing documents and communication with Chinese customers.
            \end{itemize}
            \bigskip
            % \divider
            
            \cvevent{Flight Control Engineer }{| Autel Robotics Europe GmbH}{06.2016 -- 01.2020}{Ismaning, Germany}
            \begin{itemize}
                \item Participate in development of flight control system for a 5kg tilt-rotor eVTOL UAV.
                \item Focus on the attitude loop design and control allocation for the tilting mechanism.
                \item Organize flight tests and pilot the aircraft with aid of autopilots.
                \item Responsible for organizing administrative communication and reports to the Chinese headquarter of the project.
                % \item Provide both technical and administrative supports to other projects.
            \end{itemize}
        % ----- EXPERIENCE -----
        
        % ----- EDUCATION -----
        \cvsection{Education}
            \cvevent{Doctoral Candidate }{| Technical University of Munich (TUM)}{06.2016 -- Now (Expected Grad.: 2022)}{Munich, Germany}
            \begin{itemize}
                \item \textbf{Research Direction}: \small Redundant System Control Allocation Optimization \& Protection
            \end{itemize}
            % \smallskip
            % \divider
            
            \cvevent{M.Sc. Aerospace Engineering }{| Nanyang Technological University (NTU)}{08.2013 -- 02.2016}{Singapore}
            \begin{itemize}
                \item \textbf{CGPA}: 4.92 \small (CAP: 1.06) \normalsize  \hspace{1em} \textbf{Major Subjects}: Flight Dynamics and Control, Light Weight Structures
            \end{itemize}
            % \smallskip
            % \divider
            
            % \cvevent{Uncompleted}{| Beihang University}{09.2012 -- 07.2013}{Beijing, China}
            % \begin{itemize}
            %     \item First year courses complete.
            %     \item 50\% finished towards M.Sc. Aerospace Engineering.
            % \end{itemize}
            % \divider
            
            \cvevent{B.Eng. Aircraft Design Engineering }{| Beijing Institute of Technology (BIT)}{08.2008 -- 07.2012}{Beijing, China}
            \begin{itemize}
                \item \textbf{GPA}: 85/100 \hspace{5em} \textbf{Major Subjects}: Flight Mechanics, Aerodynamics
            \end{itemize}
        % ----- EDUCATION -----
        
        % ----- PROJECTS -----
        % \cvsection{Projects}
        %     \cvevent{V280 }{\cvrepo{| \faGithub}{https://github.com/user/repo}\cvrepo{| \faGlobe}{https://repo-demo.com/}}{02.2020 -- Now}{}
        %     \begin{itemize}
        %         \item Item 1
        %         \item Item 2
        %     \end{itemize}
        %     \divider
        
        
        \renewcommand{\cvevent}[4]{%
          {\large\color{emphasis}#1}
          \ifstrequal{#2}{}{}{\large{\color{accent}#2}}
          \par\medskip\normalsize
          \ifstrequal{#3}{}{}{{\small\makebox[0.5\linewidth][l]{\color{accent}\faCalendar\color{emphasis}~#3}}}
          \ifstrequal{#4}{}{}{{\small\makebox[0.5\linewidth][l]{\color{accent}\faDev\color{emphasis}~#4}}}\par
          \medskip\normalsize%\faMapMarker
        }
        
        \cvsection{Projects}
            \cvevent{V280 - 280kg eVTOL with Hybrid Configuration }{| TUM-FSD}{02.2020 -- Now}{Industry Project}
            \begin{itemize} \justifying
                \item \textbf{Role}: Lead of control allocation development
                % \item Major Requirements:
                % % \vspace{1mm}
                % \begin{itemize}
                % \setlength\itemsep{0.1ex}
                \item \textbf{Main Contribution}: Develop the control allocation algorithm for the complete flight envelope (hover, transition and cruise flight). In addition to nominal operation, the algorithm can safely operate with a critical hover propeller failure given information from fault detection and monitoring scheme.
                \item \textbf{Development}: Strictly under code-compliant guidelines specified by DO-331 MBD design process (standard libraries, Model Advisor checks, Simulink Test Harnesses, and guaranteed WCET on target platform); automatic tests by continuous integration at push time \& issue tracking with GitLab; requirements tracking in Polarion.
                % \item \textbf{Status}: The implementation is finished over the whole flight envelope and being verified in the hardware-in-the-loop environment.
                % \item Additionally build an AMS opitmization framework to support system controllability analysis.
                % \end{itemize}
            \end{itemize}
            \bigskip
            % \divider

            \cvevent{Extended F-16 Modelling and Control}{| TUM-FSD}{05.2015 -- 07.2017}{Research Project}
            \begin{itemize} \justifying
                \item \textbf{Role}: Lead of simulation model and baseline control design
                \item \textbf{Main Contribution}: Derive \& build the simulation model and closed-loop framework. Verified the INDI concept and compared several different control allocation schemes.
                \item \textbf{Objective}: The highly redundant and nonlinear model serves as a conceptualization platform before a strategy is utilized to industry projects.
                \item \textbf{Development}: The closed-loop framework is setup in  MATLAB/Simulink with FlightGear visualization. Implementation details are found in Publications [1][3].
            \end{itemize}
            \bigskip
            % \divider

            \cvevent{Dragonfish - 5kg Tilt-rotor Transition UAV}{| Autel Robotics Europe GmbH}{06.2016 -- 01.2020}{Industry Project}
            \begin{itemize} \justifying
                \item \textbf{Role}: Lead of control allocation development
                \item \textbf{Main Contribution}: Developed the control allocation scheme for the entire envelope. Deeply involved in the autopilot design of attitude and velocity control, debugging and tuning the closed-loop behavior. The controller provides a smooth transition and unified behavior between hover and wingborne flight.
                \item \textbf{Development}: Deploys MATLAB/Simulink and code generation tool chain with some hundred hours of on-site flight test and tuning. By spring 2020, the complete autopilot software is matured for product pre-release.
                \item In addition, extensively contribute to administrative communication processes in the cooperation between the university team and the Chinese headquarter.
            \end{itemize}
            \bigskip
            % \divider
            
            \cvevent{Manta Ray - 30kg eVTOL Delivery Drone}{| Phoenix-Wings GmbH}{06.2020 -- Now}{Industry Project}
            \begin{itemize} \justifying
                \item \textbf{Role}: Programmer of Linearization Tool in C++
                \item \textbf{Main Contribution}: Design and implement the linearization tool in C++, as well as the MATLAB/Simulink interface for model verification and further design tasks.
                \item \textbf{Development}: Use generated C code of the Simulink model as interface. Code in Linux OS with open source Eclipse IDE and CMake projects. Use Ipopt and lapack library as trim solver.
                \item Additionally participate in sensor calibration processes and extensive software and on-site flight testing for the product tool chain; support in administrative processes with Chinese customers.
            \end{itemize}
        % ----- PROJECTS -----
        
        % ----- PROJECTS -----
        \cvsection{First Author Publications}
        \begin{enumerate} \justifying \small 
            \item \emph{Attainable Moment Set Optimization to Support Configuration Design: A Required Moment Set Based Approach}
            \item \emph{Saturation Protection with Pseudo Control Hedging: A Control Allocation Perspective}
            \item \emph{Control Allocation Framework with SVD-based Protection for a Tilt-rotor VTOL Transition Air Vehicle}
            \item \emph{Modeling and Incremental Nonlinear Dynamic Inversion Control for a Highly Redundant Flight System}
            \item \emph{Control Allocation Framework for a Tilt-rotor Vertical Take-off and Landing Transition Aircraft}
        \end{enumerate}
        
        
    \end{paracol}
\end{document}
